\section{Deleting one member of the simultaneous integer family}
\label{sec:delete-integer-family-member}

The new CRT bridge needs the integer quotient $j$, now interpreted as
$J$, but no longer needs the old integer $o$. The general and specified
rational interfaces can therefore be instantiated with eleven input
blocks. This is a new application of their finite-family theorem; it
does not assert that a rational zero of an eleven-block system extends
to a zero of the former twelve-block system with a prescribed $o$.

Here is the required family argument. Keep the blocks indexed by
\[
              0,1,2,3,4,5,6,7,8,10,11,
\]
with the original ordered labels $b,g,t,m,w,s,h,c,\gamma,j,i$.
The reference block $b$, common point roles, common input numerator,
common contact gauge, and reciprocal $i$ chart remain. For $e=2$
general contacts or $e=3$ specified contacts, define the gauge by the
sum of the retained weights only. Thus its number of summands is now
$m=22$ or $m=33$, respectively.

The soundness proof of the global gauge uses a consumed prime of one
retained weight, its own complementary sum, the unchanged scalar and
point relations, and its own pole obstruction. It makes no assumption
that another input block exists. All remaining height estimates use
the same finite-family inequalities, with fewer summands; the former
twelve-block upper bounds are still sufficient. Hence every remaining
distinguished input is an integer under the same arithmetic premises.
The affine contact chart is reconstructed with the seven common
coefficient registers exactly as before.

For completeness, start with the eleven desired integer inputs and
choose their point data and aggregate shifts by the existing ordered
simultaneous preparation. Take the common contact spacing sufficiently
large to separate the translated finite lists. The displayed indices
are distinct integers, so the interval-separation proof applies even
though index $9$ has been omitted. Invoke the same finite-family prime
supplier on their union, with the same finite avoidance and all-sign
common-index congruences. At each selected contact, only its own weight
has negative odd valuation, and all other retained weights are integral.
At foreign contacts, the same rational-square parity argument applies.
At the prepared finite places and at infinity,
\[
 q_\nu=\kappa\lambda^{-2}+O(\lambda^{-3}),\qquad
 \rho=m\kappa\lambda^{-2}+O(\lambda^{-3}),\qquad
 \frac{(\rho-q_\nu)^2}{q_\nu}
       =(m-1)^2\kappa\lambda^{-2}+O(\lambda^{-3})\longrightarrow0.
\]
No division by $m$ or $m-1$ is needed. The physical tests are bounded
on the previously chosen neighborhoods; their nonconstant detector
terms tend to zero. The actual pole arguments thus retain the required
local square classes at $2,3,17$ and the same real limit. Choosing
sufficient finite divisibility and a large real slope gives every
remaining native-pole witness. This proves the eleven-block
simultaneous completeness assertion.

These are precisely the arguments in
\emph{One contact gauge in a combined quadratic system},
\nolinkurl{undecidability/papers/global_gauge_quartic_checkpoint.tex},
Sections 2 and 3, and
\emph{A common arithmetic progression of contact arguments},
\nolinkurl{undecidability/general_contact_v94/affine_contact_progressions.tex},
with a smaller finite family. No new prime-supplier premise, positivity
assumption on an arbitrary rational zero, or pairwise radical
disjointness assumption is introduced.

\paragraph{Complete coefficient audit.}
In the signed affine-contact tables, the literal private coordinates
are $o$ and every name beginning \texttt{int\_o\_}. There are $61$
of them in the general table and $76$ in the specified table. The
private integer-certificate rows number $46$ and $59$. Each interface
also has four private \texttt{global.o.} input-pair rows, one private
\texttt{scalar\_norm.o.} row, and one private
\texttt{detector\_trace.o.} row. Their deletion removes $52$ and
$65$ quadratic rows, respectively.

Apart from the old source row $I-d-oF$, the only occurrences of these
private coordinates outside those rows lie in three rows of the last
reciprocal-$i$ contact. They are exactly the eliminated global-sum
contributions of the removed weights. The general contributions are
\texttt{int\_o\_153} and
\texttt{int\_o\_last\_contact\_weight}; the specified contributions
are \texttt{int\_o\_154}, \texttt{int\_o\_155}, and
\texttt{int\_o\_last\_contact\_weight}.
Deleting their terms from these three rows compiles the smaller global
sum. All other retained rows, including the common gauge guards and
the private guards of every retained block, are literally unchanged.

The CRT source then drops $I-d-oF$, replaces the old definition of $H$,
and adds one coordinate $CF$ with the equation $CF-cF=0$. Its new
quadratic $H$ equation is
\[
              H-c-2rF_2-F_2+cF_2-(CF)j=0,
              \qquad F_2=F^2.
\]
Thus the source row count is unchanged, and its additional coordinate
offsets one of the deleted private coordinates. The full counts become
$1043-61+1=983$ variables and $907-52=855$ rows in the general
interface, and $1223-76+1=1148$ variables and $1063-65=998$ rows in
the specified interface. The exact coefficient checker independently
verifies every private support, every deleted global-sum term, every
unchanged row, the new source equations, coefficient domains, guards,
degree, counts, and complete coordinate support.
